\documentclass{situationist-preprint}

\title{The Toxin Puzzle}
\msid{ACT-06-b}
\msfamily{ACT-06}
\msdiscipline{Action, Agency and Free Will}
\mstype{Companion manuscript}
\mskeywords{toxin puzzle; intention; commitment; reasons for intending; planning theory; Kavka}

\begin{document}
\maketitle

\begin{abstract}
Kavka's toxin puzzle offers a large reward for intending at midnight to drink a mildly unpleasant toxin the next afternoon; the reward is paid on the intention alone, whether or not one drinks. It seems that one cannot form the intention, though one has excellent reason to have it. We argue that the account of intention as committed planning in ACT-06 explains why. Commitment settles a practical question, and a question about what to do is settled only by reasons bearing on doing it. The reward is a reason to have the intention, not to drink, so it cannot settle the question. The puzzle also confirms that intention does not reduce to belief and desire: the agent may want to intend and expect to drink without intending.
\end{abstract}

\section{The puzzle}

A billionaire offers you a million dollars if, at midnight tonight, you intend to drink a toxin tomorrow afternoon that will make you ill for a day \citep{Kavka1983}. A reliable brain scanner will detect the intention, and the money will be in your account by morning, whether or not you then drink. You have overwhelming reason to intend to drink. But you know that tomorrow you will have no reason to drink, and every reason not to. It seems you cannot form the intention, and that if you could, it would be irrational to carry it out.

\section{Commitment and settling}

ACT-06-a defines an intention as an attitude toward a plan that involves commitment (Definition~2.3), and commitment as settling, a disposition to treat the question whether to execute the plan as closed, together with constraint on further reasoning (Definition~2.2).

\begin{proposition}\label{p:settle}
A reason that bears only on having an intention, and not on executing its plan, cannot settle the question whether to execute the plan.
\end{proposition}

\begin{proof}
To settle a practical question is to close it in a way that the agent is disposed not to reopen without new information or a change in her reasons (Definition~2.2(i) of ACT-06-a). The question is whether to drink. A reason that bears only on having the intention does not favor drinking over not drinking, and so cannot close the question in favor of drinking. If the agent nevertheless treated the question as closed, she would be disposed not to reconsider it, although the only considerations bearing on it, which she already foresees, favor not drinking. That disposition would not be the settling of the question by her reasons, but the refusal to consider them.
\end{proof}

The proposition explains Kavka's verdict without the claim that intentions are beliefs that one will act. The agent cannot form the intention directly, because what would settle the question is absent. She may form it indirectly, by arranging that she will have reason to drink, for instance by a side bet with a friend; then the intention is settled by reasons for drinking, and executing it is rational. The constraint component of commitment confirms the diagnosis: drinking is not a means to the reward, which is paid at midnight, so the means--end requirement of Proposition~3.4 of ACT-06-a gives the plan of drinking no support from the end of being rewarded.

\paragraph{Non-reduction.} The puzzle supplies a natural instance of Proposition~3.2 of ACT-06-a. An agent who wants the reward wants very much to intend to drink, and, if she knows herself to be compliant, may believe that she will drink. She may still not have intended, since the question whether to drink has not been settled. Strong desire for the intention and the expectation of acting fall short of the commitment intention involves.

\section{An objection}

\begin{objection}[Rational plans]
On a planning view, it is rational to adopt and carry out a plan whose adoption makes one better off, and the plan to drink does \citep{Gauthier1994}. So the agent can intend, and should drink.
\end{objection}

\begin{reply}
The view on which plans are assessed as wholes is the strongest reply to Kavka, but it does not cohere with the stability of intention. An intention resists reconsideration because reconsidering is costly and because its grounds are expected to hold; neither applies when the grounds of the plan are known in advance to lapse at the moment of execution. \citet{Bratman1987} defends this limit on the stability of intentions, and our account inherits it: commitment is a disposition not to reopen a question without new reason, not a disposition to act against the reasons that one foresaw when one closed it.
\end{reply}

\section{Conclusion}

The toxin puzzle shows that one cannot intend at will for the sake of a reward for intending, and the account of intention as committed planning explains why: commitment settles a question, and only reasons bearing on its answer can settle it. The puzzle thereby supports the conclusion of ACT-06, that intentions are plans with commitment and do not reduce to beliefs and desires, which involve no commitment.

\begin{thebibliography}{9}
\bibitem[Bratman(1987)]{Bratman1987} Bratman, M.~E. (1987). \emph{Intention, Plans, and Practical Reason}. Harvard University Press.
\bibitem[Gauthier(1994)]{Gauthier1994} Gauthier, D. (1994). Assure and threaten. \emph{Ethics}, 104(4), 690--721.
\bibitem[Kavka(1983)]{Kavka1983} Kavka, G.~S. (1983). The toxin puzzle. \emph{Analysis}, 43(1), 33--36.
\end{thebibliography}

\end{document}
